\section{Historia antes de los LLM Agent}

\subsection{Text Agent basados en reglas}

\begin{knowledgebox}{ELIZA (1966)}
Uno de los primeros chatbots, capaz de dar la impresión de «inteligencia» usando solo unas cuantas reglas (preguntar insistentemente o repetir lo que decía el usuario). Sin embargo, el método basado en reglas tiene limitaciones serias:
\begin{itemize}
    \item Altamente específico de dominio (domain-specific)
    \item Cada dominio nuevo requiere volver a escribir las reglas
    \item No se puede escalar a dominios abiertos complejos
\end{itemize}
\end{knowledgebox}

\subsection{Text Agent basados en aprendizaje por refuerzo}

Antes de la aparición de los LLM, otro paradigma popular era construir Text Agent con RL (por ejemplo, agentes para juegos de texto):

\begin{itemize}
    \item Tratar las observaciones y acciones de texto como los píxeles y las entradas de teclado en un videojuego
    \item Usar RL para optimizar la señal de recompensa
    \item \textbf{Limitación}: específico de dominio, requiere una señal de recompensa precisa, costo de entrenamiento alto
\end{itemize}

\subsection{El cambio revolucionario de los LLM}

Los LLM, entrenados únicamente con next-token prediction sobre enormes cantidades de texto, pueden resolver diversas tareas nuevas en el momento de la inferencia mediante prompting. Esta \textbf{generalidad} y esta \textbf{capacidad de aprendizaje con pocos ejemplos} los convierten en el «cerebro» ideal para construir Agent.

\subsection{Resumen de la sección}

Del symbolic AI agent basado en reglas al deep agent basado en RL, y de ahí al reasoning agent basado en LLM, la representación interna del Agent pasó de la lógica simbólica a los embeddings neuronales y de ahí al lenguaje natural, avanzando paso a paso hacia una mayor flexibilidad y generalidad.
